% ═══════════════════════════════════════════════════════════
% ── TEMPLATE SLIDES ───────────────────────────────────────
% ═══════════════════════════════════════════════════════════

\begin{frame}[plain]
  \vfill
  \begin{center}
    {\LARGE\color{slate900}\textbf{Template Slides}}\\[8pt]
    {\color{slate700}Reference examples for building new slides}
  \end{center}
  \vfill
\end{frame}

% ── Template: Section Divider ───────────────────────────────
\begin{frame}[plain]
  \vfill
  \begin{center}
    {\LARGE\color{slate900}\textbf{Section Title}}\\[8pt]
    {\color{slate700}Brief subtitle or context}
  \end{center}
  \vfill
\end{frame}

% ── Template: Content Slide ─────────────────────────────────
\begin{frame}{Content Slide Example}
  Key point or framing sentence at the top.

  \begin{itemize}
    \item First point with enough detail to stand alone
    \item Second point
      \begin{itemize}
        \item Supporting detail
        \item Supporting detail
      \end{itemize}
    \item Third point
  \end{itemize}
\end{frame}
\note{Example note: mention the 16,000+ Robotarium experiments here. Pause before transitioning.}

% ── Template: Equation Slide ────────────────────────────────
\begin{frame}{Equation Slide Example}
  The safety filter solves a quadratic program at each timestep:

  \vskip10pt
  \begin{block}{Active Safety Filter}
    \vskip4pt
    \begin{equation*}
      \min_{u} \quad \| u - u_{\text{des}} \|^2
    \end{equation*}
    \vskip-4pt
    \begin{equation*}
      \text{s.t.} \quad L_f h(x) + L_g h(x)\, u \geq -\alpha\big(h(x)\big)
    \end{equation*}
    \vskip2pt
  \end{block}

  \vskip8pt
  \begin{itemize}
    \item $h(x) \geq 0$ defines the safe set
    \item Minimally invasive: follows $u_{\text{des}}$ when safe
    \item Small QP, solves in real time
  \end{itemize}
\end{frame}
\note{Walk through each term slowly. Dennis will know this cold. Peter will appreciate the connection to his trusted autonomy work. If asked about scalability, be honest about high-dimensional challenges.}

% ── Template: Two-Column Slide ──────────────────────────────
\begin{frame}{Two-Column Slide Example}
  \begin{columns}[T]
    \begin{column}{0.48\textwidth}
      \textbf{Left: Text or Math}

      \vskip8pt
      \begin{itemize}
        \item Explain the concept
        \item Key equations or points
        \item Connection to Archer
      \end{itemize}
    \end{column}
    \begin{column}{0.48\textwidth}
      \textbf{Right: Figure}

      \vskip8pt
      % \includegraphics[width=\textwidth]{figures/placeholder.png}
      \begin{center}
        {\color{slate400}\framebox(180,100){Figure goes here}}
      \end{center}
    \end{column}
  \end{columns}
\end{frame}

% ── Template: Image Slide ───────────────────────────────────
\begin{frame}{Full Image Slide Example}
  \begin{center}
    % \includegraphics[width=0.85\textwidth]{figures/placeholder.png}
    {\color{slate400}\framebox(400,180){Full-width figure or screenshot}}
  \end{center}

  \vskip4pt
  Brief caption or takeaway below the image.
\end{frame}

% ── Template: Side-by-Side Images ─────────────────────────
% Pre-crop images to same aspect ratio (e.g., 4:3) with Python/PIL.
% Equal columns + equal ratios = matched heights automatically.
% Use \vspace{5mm} before source to visually center between title rule and footnote.
\begin{frame}{Side-by-Side Image Example}
  \vfill
  \begin{columns}[c]
    \begin{column}{0.495\textwidth}
      % \includegraphics[width=\textwidth]{figures/left_crop.png}
      {\color{slate400}\framebox(210,158){Left image (4:3)}}
    \end{column}
    \begin{column}{0.495\textwidth}
      % \includegraphics[width=\textwidth]{figures/right_crop.png}
      {\color{slate400}\framebox(210,158){Right image (4:3)}}
    \end{column}
  \end{columns}
  \vfill
  \source{Source citation here}
\end{frame}
